% ============================================================================
% 元数据配置 - config.tex
% 包含论文标题、作者信息、摘要、关键词等元数据
% ============================================================================

% --- 论文标题 ---
\newcommand{\papertitle}{顾及风险时空异质性的城市低空无人机航路规划}

% --- 作者信息 ---
\newcommand{\authorname}{xxx, xxx, xxx}
\newcommand{\authoremail}{xxx@example.com}
\newcommand{\authoraddress}{xxxx大学xxx院，城市邮编}
\newcommand{\correspondingnote}{通讯作者}

% --- 摘要内容 ---
\newcommand{\abstracttext}{
无人机（UAV）在城市环境中日益广泛地应用于物流配送、巡检和应急响应等领域，但其运行对地面行人、车辆及建筑物构成坠机撞击、财产损失与噪声滋扰等第三方风险。现有风险感知路径规划方法多假设静态环境与恒定坠机概率，难以刻画城市微气象、建筑峡谷效应及人口昼夜潮汐所导致的风险时空异质性。

为此，本文构建一套多维第三方风险评估模型与时间依赖路径规划算法相结合的一体化框架。在风险建模层面，基于比例风险理论融合风场、降雨与城市峡谷因子建立时变坠机概率，结合POI引力模型与昼夜激活函数生成时空动态人口及车辆暴露密度，并从致死、财产损失与噪声影响三个维度量化风险代价。在求解算法层面，提出时间依赖风险感知A*算法（TD-RiskA*），将时间维从搜索图拓扑中剥离并作为标签状态沿路径传播，使每次节点扩展可在线查询动态风险张量；引入累积危险率加法传播机制以保障长距离搜索的数值稳定性，并设计多标签支配剪枝框架；理论分析与实验均表明，在本文单调代价结构下单标签覆盖即达最优，计算开销与传统三维 A* 同阶。

合成场景与真实城市场景实验表明，所提方法在极小绕行代价下显著降低核心城区的物理撞击风险与敏感时段噪声滋扰。相较于静态启发式规划，TD-RiskA*可依据城市时间节律与多目标权重偏好生成时空自适应的安全飞行路径。
}

% --- 关键词 ---
\newcommand{\paperkeywords}{
第三方风险 \sep 时空异质性 \sep 无人机 \sep 低空经济 \sep 路径规划
}

% --- MSC 分类代码 ---
\newcommand{\msccodes}{90Cxx \sep 68Txx}